\section{Latihan}

\begin{latihan}[Analisis Game Favorit]
Pilih satu game favorit Anda dan analisis menggunakan MDA framework:
\begin{enumerate}
\item Identifikasi minimal 5 mechanics utama
\item Jelaskan dynamics yang muncul dari mechanics tersebut
\item Identifikasi aesthetics yang dirasakan saat bermain
\item Bagaimana mechanics, dynamics, dan aesthetics saling berhubungan?
\end{enumerate}
\end{latihan}

\begin{latihan}[Analisis Genre]
Pilih 3 game berbeda dari genre yang berbeda dan identifikasi:
\begin{enumerate}
\item Genre utama dan sub-genre (jika ada)
\item Karakteristik gameplay yang membuat game tersebut termasuk dalam genre tersebut
\item Elemen dari genre lain yang mungkin ada dalam game tersebut
\end{enumerate}
\end{latihan}

\begin{latihan}[Membuat GDD]
Buat Game Design Document untuk konsep game sederhana dengan mencakup minimal:
\begin{enumerate}
\item Executive Summary
\item Game Overview
\item Core Gameplay Mechanics
\item Art Style
\item Development Timeline (estimasi)
\end{enumerate}
\end{latihan}

\begin{latihan}[Membuat Development Timeline]
Buat timeline pengembangan untuk game sederhana dengan:
\begin{enumerate}
\item Breakdown fase pre-production, production, post-production
\item Estimasi waktu untuk setiap fase
\item Identifikasi milestone utama
\item Identifikasi risiko potensial dan mitigasinya
\end{enumerate}
\end{latihan}

\begin{latihan}[Merancang Tim]
Untuk game sederhana (misalnya: 2D platformer), rancang struktur tim minimal dengan:
\begin{enumerate}
\item Identifikasi peran yang diperlukan
\item Deskripsi tanggung jawab setiap peran
\item Estimasi jumlah personel untuk setiap peran
\item Identifikasi tools yang diperlukan untuk kolaborasi
\end{enumerate}
\end{latihan}
